\section{Oeljeklaus--Toma manifolds}
\label{sec:oeljeklaus-toma}
\index{Oeljeklaus--Toma manifold}\index{Inoue surface}

The compact manifolds studied so far in this chapter carry many closed positive currents, coming from divisors, fibrations or blowups. Inoue surfaces and their higher-dimensional analogues, the Oeljeklaus--Toma manifolds, behave very differently. They are built from number fields, contain few or no complex subvarieties, and carry a natural closed semipositive $(1,1)$-form whose kernel is a holomorphic foliation by complex affine spaces. These leaves are Liouville and fill out real tori. As a result, every potential in this class is a continuous convex function of the transverse coordinates, so the class has no singular potentials at all, even though it contains positive currents supported on real hypersurfaces. In dimension three, a natural Gauduchon metric has an exactly quadratic mass functional of indefinite sign. Its masses lie on both sides of the background mass, and whether they remain bounded is an open question (\cref{qu:ot-bounded-mass}).

\subsection{The construction}

Let $K$ be a number field of degree $s+2t$, with $s\geq1$ real embeddings $\sigma_1,\ldots,\sigma_s$ and $t\geq1$ pairs of complex conjugate embeddings $\sigma_{s+k},\overline{\sigma}_{s+k}$, $1\leq k\leq t$. Write $\mathcal O_K$ for its ring of integers and $\mathcal O_K^{*,+}$ for the group of totally positive units, that is, of units $\varepsilon$ with $\sigma_j(\varepsilon)>0$ for $1\leq j\leq s$. Let $\mathbb H\subset\mathbb C$ be the upper half-plane. Write points of $\mathbb H^s\times\mathbb C^t$ as $(w,z)$, with $w_j=x_j+\mathrm i y_j$, and put
\[
 \sigma_{\mathbb R}(a)=(\sigma_1(a),\ldots,\sigma_s(a)),\qquad
 \sigma_{\mathbb C}(a)=(\sigma_{s+1}(a),\ldots,\sigma_{s+t}(a)).
\]
An element $a\in\mathcal O_K$ acts by the translation $(w,z)\mapsto(w+\sigma_{\mathbb R}(a),z+\sigma_{\mathbb C}(a))$, and $\varepsilon\in\mathcal O_K^{*,+}$ acts by
\[
 (w,z)\longmapsto\bigl(\sigma_1(\varepsilon)w_1,\ldots,\sigma_s(\varepsilon)w_s,\,
       \sigma_{s+1}(\varepsilon)z_1,\ldots,\sigma_{s+t}(\varepsilon)z_t\bigr).
\]
A subgroup $U\subset\mathcal O_K^{*,+}$ is \emph{admissible} if its image
\[
 \Lambda_U=\{(\log\sigma_1(\varepsilon),\ldots,\log\sigma_s(\varepsilon)):\varepsilon\in U\}
\]
is a lattice in $\mathbb R^s$. By Dirichlet's unit theorem, when $t=1$ every subgroup of finite index in $\mathcal O_K^{*,+}$ is admissible. For admissible $U$, Oeljeklaus and Toma showed that the semidirect product $\Gamma=\mathcal O_K\rtimes U$ acts freely and properly discontinuously on $\mathbb H^s\times\mathbb C^t$ with compact quotient. The compact complex manifold
\[
 X=X(K,U)=(\mathbb H^s\times\mathbb C^t)/\Gamma,
 \qquad n=\dim X=s+t,
\]
is the \emph{Oeljeklaus--Toma manifold} of type $(s,t)$.\index{Oeljeklaus--Toma manifold!construction} For $s=t=1$ these are exactly the Inoue surfaces $S_M$. The sources are discussed in \cref{sec:history-closed-loci}.

The functions $y_j=\operatorname{Im}w_j$ are invariant under the translations and are multiplied by $\sigma_j(\varepsilon)$ under the unit $\varepsilon$. Hence the closed real $(1,1)$-form
\[
 \theta=-\ddc\sum_{j=1}^s\log y_j
       =\sum_{j=1}^s\frac{\mathrm i\,\mathrm dw_j\wedge\mathrm d\overline w_j}{8\pi y_j^2}
\]
on $\mathbb H^s\times\mathbb C^t$ is $\Gamma$-invariant and descends to a closed semipositive form on $X$, still denoted $\theta$.\index{Oeljeklaus--Toma manifold!transverse form} It has rank $s$ everywhere, and its kernel is spanned by the vectors $\partial/\partial z_k$. The potential $-\sum_j\log y_j$ itself is not invariant, since a unit adds $-\sum_j\log\sigma_j(\varepsilon)$ to it.

\subsection{Potentials of the transverse form}

Two properties of the action force rigidity: the leaves $\{w\}\times\mathbb C^t$ are Liouville, and the real parts of the translations are dense.

\begin{lemma}\label{lem:ot-real-density}
Since $t\geq1$, the subgroup $\sigma_{\mathbb R}(\mathcal O_K)$ is dense in $\mathbb R^s$.
\end{lemma}
\begin{proof}
A subgroup $G$ of $\mathbb R^s$ is dense unless there is a nonzero linear form $\xi$ on $\mathbb R^s$ with $\xi(G)\subseteq\mathbb Z$: the closure of $G$ is the sum of a linear subspace $V$ and a discrete subgroup of a complement of $V$, and if $V\neq\mathbb R^s$ one can take $\xi$ vanishing on $V$ and integral on that discrete subgroup. Suppose $\xi(\sigma_{\mathbb R}(a))=\sum_j\xi_j\sigma_j(a)\in\mathbb Z$ for every $a\in\mathcal O_K$. Since $\mathcal O_K$ spans $K$ over $\mathbb Q$, the map $a\mapsto\sum_j\xi_j\sigma_j(a)$ is a $\mathbb Q$-linear map $K\to\mathbb Q$. The trace form of $K$ is nondegenerate, so there is $b\in K$ with
\[
 \sum_{j=1}^s\xi_j\sigma_j(a)=\operatorname{Tr}_{K/\mathbb Q}(ba)
 =\sum_{\sigma}\sigma(b)\sigma(a)
 \qquad(a\in K),
\]
where $\sigma$ runs over all $s+2t$ embeddings of $K$ into $\mathbb C$. Distinct field embeddings are linearly independent over $\mathbb C$, by Dedekind's lemma on characters. Comparing the coefficients of a complex embedding $\sigma_{s+1}$, which occurs only on the right, gives $\sigma_{s+1}(b)=0$. Hence $b=0$, and then $\xi_j=\sigma_j(b)=0$ for every $j$.
\end{proof}

For $\varepsilon\in U$ and $y\in(0,\infty)^s$, write $\sigma(\varepsilon)y=(\sigma_1(\varepsilon)y_1,\ldots,\sigma_s(\varepsilon)y_s)$.

\Needspace{12\baselineskip}
\begin{theorem}\label{thm:ot-transverse-potentials}\index{Oeljeklaus--Toma manifold!transverse potentials}\index{Plurisubharmonic function!rigidity along a foliation}
Let $\mathcal F_U$ be the set of convex functions $f:(0,\infty)^s\to\mathbb R$ such that
\begin{equation}\label{eq:ot-automorphy}
 f(\sigma(\varepsilon)y)=f(y)-\sum_{j=1}^s\log\sigma_j(\varepsilon)
 \qquad(\varepsilon\in U).
\end{equation}
For $f\in\mathcal F_U$, the function $u_f(w,z)=f(y)+\sum_j\log y_j$ is $\Gamma$-invariant, and $f\mapsto u_f$ is a bijection from $\mathcal F_U$ onto $\PSH(X,\theta)$. In particular, every $\theta$-plurisubharmonic function on $X$ is continuous and depends only on $y=\operatorname{Im}w$.
\end{theorem}
\begin{proof}
We first note that a function $g$ on $\mathbb H^s$ that depends only on $y$ is plurisubharmonic if and only if, as a function on $(0,\infty)^s$, it is convex (with values in $[-\infty,+\infty)$ and not identically $-\infty$). For smooth $g$ this follows from $\partial^2g/\partial w_j\partial\overline w_k=\frac14\partial^2g/\partial y_j\partial y_k$. In general, convolution with a smooth radial kernel in $w$ produces smooth functions of $y$ alone; they decrease to $g$ if $g$ is plurisubharmonic, and they converge locally uniformly to $g$ if $g$ is convex, because a convex function that is not identically $-\infty$ on the open convex set $(0,\infty)^s$ is finite and continuous there. Since decreasing limits of convex functions are convex and locally uniform limits of plurisubharmonic functions are plurisubharmonic, the claim follows.

The identity \eqref{eq:ot-automorphy} says exactly that $u_f$ is invariant under $U$, and $u_f$ is clearly invariant under the translations. On the universal cover, $\theta+\ddc u_f=\ddc f(y)\geq0$. Hence $u_f\in\PSH(X,\theta)$, and $f$ is recovered from $u_f$, so the map is injective.

Conversely, let $u\in\PSH(X,\theta)$, and let $\tilde u$ be its lift to $\mathbb H^s\times\mathbb C^t$. Since $\theta+\ddc\tilde u=\ddc(\tilde u-\sum_j\log y_j)$ there, the function $q=\tilde u-\sum_j\log y_j$ is plurisubharmonic on $\mathbb H^s\times\mathbb C^t$.

\emph{Step 1: constancy on the leaves.} Fix $w^0\in\mathbb H^s$. The restriction of $q$ to the slice $\{w^0\}\times\mathbb C^t$ is plurisubharmonic or identically $-\infty$, and it is bounded above because $u$ is upper semicontinuous on the compact manifold $X$. A plurisubharmonic function on $\mathbb C^t$ that is bounded above is constant, since its restriction to every complex line through two given points is a subharmonic function on $\mathbb C$ bounded above. Hence $\tilde u(w,z)=\tilde u(w,0)$ for all $(w,z)$. In particular, the function $F(w)=q(w,0)$ is not identically $-\infty$, because $\tilde u$ is not, and it is plurisubharmonic on $\mathbb H^s$.

\emph{Step 2: invariance under real translations.} The translations by $\mathcal O_K$ and Step~1 give $F(w+\sigma_{\mathbb R}(a))=F(w)$ for every $a\in\mathcal O_K$, since $\log y_j$ does not change under real translations. For $r>0$, the averages $F_r$ of $F$ over the Euclidean balls of radius $r$ are defined on $\{w:\min_jy_j>r\}$ and continuous there, because $F$ is locally integrable, and they are invariant under the same translations. By \cref{lem:ot-real-density} and continuity, each $F_r$ is invariant under all translations by $\mathbb R^s$. Since $F_r$ decreases to $F$ as $r\downarrow0$, the function $F$ is invariant under all real translations as well, so $F(w)=f(y)$ for a function $f$ on $(0,\infty)^s$.

\emph{Step 3: convexity and automorphy.} Since $F=f(y)$ is plurisubharmonic, the first paragraph of the proof shows that $f$ is convex, finite and continuous. Finally, $u=u_f$, and invariance of $u$ under $U$ is \eqref{eq:ot-automorphy}.
\end{proof}

\begin{corollary}\label{cor:ot-no-singular-potentials}
On every Oeljeklaus--Toma manifold, all functions in $\PSH(X,\theta)$ are continuous. Consequently, they all have the ordinary singularity type of $0$, all their Lelong numbers vanish, all their multiplier ideals are trivial, and $P_\theta[u]=V_\theta=0$ for every $u\in\PSH(X,\theta)$. Moreover $\mathrm{MA}_\theta(u)=0$ for every $u$, so $\theta$ is nef but not big.
\end{corollary}
\begin{proof}
Continuous functions on a compact manifold are bounded, so any two of them are equivalent, and their Lelong numbers and multiplier ideals are those of a bounded function. Since $\theta\geq0$, the constant $0$ belongs to $\PSH(X,\theta)$, so $V_\theta=0$, and $P_\theta[u]=P_\theta[0]=V_\theta$. By \cref{thm:ot-transverse-potentials}, locally $\theta+\ddc u=\ddc f(y)$ with $f$ continuous and depending on the $s$ real variables $y_j$; the Bedford--Taylor products of such a current vanish in degree $s+1\leq n$, because they are limits of the corresponding products for smooth convex approximants of $f$, whose complex Hessians have rank at most $s$. Thus $(\theta+\ddc u)^n=0$. The form $\theta$ is smooth and semipositive, hence nef. A strict barrier $\rho$ with $\theta+\ddc\rho\geq\delta\omega$ would have $\mathrm{MA}_\theta(\rho)\geq\delta^n\omega^n\neq0$, so $\theta$ is not big.
\end{proof}

Singular currents nevertheless occur in this class. Let $s=t=1$, so that $X$ is an Inoue surface and $U$ is generated by a unit $\varepsilon_0$ with $\lambda=\sigma_1(\varepsilon_0)>1$. Fix $y_0>0$, put $y_k=\lambda^ky_0$ for $k\in\mathbb Z$, and let $f$ be the continuous function that is affine on each interval $[y_k,y_{k+1}]$ with slope
\[
 a_k=-\frac{\log\lambda}{(\lambda-1)y_k},
 \qquad f(y_0)=0.
\]
Then $a_k(y_{k+1}-y_k)=-\log\lambda$, so $f(\lambda y)=f(y)-\log\lambda$, and the slopes increase because $a_{k+1}=a_k/\lambda$ and $a_k<0$. Thus $f\in\mathcal F_U$. On the universal cover,
\[
 \theta+\ddc u_f=\ddc f(y)=\frac1{4\pi}\sum_{k\in\mathbb Z}(a_k-a_{k-1})\,
       \delta(y-y_k)\,\mathrm dx\wedge\mathrm dy,
\]
with $a_k-a_{k-1}>0$. The closed positive current $\theta+\ddc u_f$ therefore has a continuous potential but is supported on the compact real hypersurface of $X$ formed by the images of the hyperplanes $\{y=y_k\}$, a three-dimensional torus that is a union of leaves of the foliation $\{w=\text{const}\}$.

On an Inoue surface, $\theta$ represents the canonical Bott--Chern class. Indeed, the Tricerri metric $\omega_T=8\pi\theta+y\,\mathrm i\,\mathrm dz\wedge\mathrm d\overline z$ has $\omega_T^2=(2/y)\,\mathrm i\,\mathrm dw\wedge\mathrm d\overline w\wedge\mathrm i\,\mathrm dz\wedge\mathrm d\overline z$, so its Chern--Ricci form $-\ddc\log(\omega_T^2/\mathrm i\,\mathrm dw\wedge\mathrm d\overline w\wedge\mathrm i\,\mathrm dz\wedge\mathrm d\overline z)$ equals $\ddc\log y=-\theta$. Hence \cref{thm:ot-transverse-potentials} describes all positive currents in $c_1^{BC}(K_X)$, and none of them has a pole.

\subsection{Monge--Amp\`ere masses on threefolds}
\index{Oeljeklaus--Toma manifold!Monge--Amp\`ere mass}

Now let $s=2$ and $t=1$, so that $X$ is a threefold. A totally positive unit has norm $\sigma_1(\varepsilon)\sigma_2(\varepsilon)|\sigma_3(\varepsilon)|^2=1$, so the form $y_1y_2\,\ddc|z|^2$ is $\Gamma$-invariant. The Hermitian form
\[
 \omega=\theta+\beta,\qquad \beta=y_1y_2\,\ddc|z|^2,
\]
descends to $X$. Write $\mathrm dV_0=(\mathrm i\,\mathrm dw_1\wedge\mathrm d\overline w_1)\wedge(\mathrm i\,\mathrm dw_2\wedge\mathrm d\overline w_2)\wedge(\mathrm i\,\mathrm dz\wedge\mathrm d\overline z)$ on the universal cover, and $\Mass_\omega(u)=\int_X(\omega+\ddc u)^3$.

Units that fix every $y_j$ are trivial: such a unit satisfies $|\sigma_3(\varepsilon)|=1$, so all its conjugates have absolute value one, it is a root of unity by Kronecker's theorem, and a totally positive root of unity in a field with a real embedding equals $1$. Hence the map $X\to T_U=\mathbb R^2/\Lambda_U$, $[w,z]\mapsto(\log y_1,\log y_2)$, has as fibers the real tori $(\mathbb R^2\times\mathbb C)/(\sigma_{\mathbb R},\sigma_{\mathbb C})(\mathcal O_K)$, and for every $\Lambda_U$-periodic function $g$ of $\varsigma=(\log y_1,\log y_2)$ there is a constant $c_X>0$, independent of $g$, with
\begin{equation}\label{eq:ot-fiber-integration}
 \int_Xg(\varsigma)\,\frac{\mathrm dV_0}{y_1y_2}
 =c_X\int_{T_U}g\,\mathrm d\varsigma.
\end{equation}

We shall use an elementary property of lattices in the plane: every lattice $\Lambda\subset\mathbb R^2$ contains vectors in each of the four open quadrants. Indeed, let $\ell>0$ be the length of a shortest nonzero vector of $\Lambda$, and let $0<\delta<\ell$. By Minkowski's theorem, the symmetric convex set $\{k:|k_1-k_2|<\delta,\ |k_1+k_2|<N\}$ contains a nonzero lattice point once $N$ is large. Such a point satisfies $(k_1+k_2)^2=2|k|^2-(k_1-k_2)^2\geq2\ell^2-\delta^2$, hence $4k_1k_2=(k_1+k_2)^2-(k_1-k_2)^2>0$; replacing $k$ by $-k$ if necessary gives a point in the open first quadrant, and the line $\mathbb R(1,-1)$ gives the other two quadrants in the same way. If $\lambda$ lies in the open first quadrant and $\mu\in\Lambda$ is independent of $\lambda$, then $N\lambda+\mu$ also lies there for large $N$, so the first quadrant contains two independent lattice vectors.

\Needspace{14\baselineskip}
\begin{proposition}\label{prop:ot-threefold-mass}
Let $X$ be an Oeljeklaus--Toma threefold of type $(2,1)$ with the Hermitian form $\omega$ above.
\begin{enumerate}
\item The form $\omega$ is Gauduchon, and $\ddc\omega=\ddc(y_1y_2)\wedge\ddc|z|^2$ is neither semipositive nor seminegative.
\item For every smooth real function $u$ on $X$,
\[
 \Mass_\omega(u)=\int_X\omega^3+\frac{3}{16\pi^3}\int_X
 \operatorname{Re}\Bigl(\frac{\partial u}{\partial w_1}\,
       \overline{\frac{\partial u}{\partial w_2}}\Bigr)\mathrm dV_0.
\]
\item Let $u\in\PSH(X,\omega)$ be bounded and depend only on $y$, and write $u=g(\varsigma)$ with $g$ periodic under $\Lambda_U$. Then $g$ is Lipschitz, $P=1-\partial_1g$ and $Q=1-\partial_2g$ are nonnegative almost everywhere, and
\[
 \Mass_\omega(u)=\frac{\int_X\omega^3}{|T_U|}\int_{T_U}PQ\,\mathrm d\varsigma.
\]
There is a constant $c_U\geq1$, depending only on $\Lambda_U$, such that every such $u$ satisfies $c_U^{-1}\int_X\omega^3\leq\Mass_\omega(u)\leq c_U\int_X\omega^3$.
\item There are smooth $u\in\PSH(X,\omega)$, depending only on $y$, with $\Mass_\omega(u)>\int_X\omega^3$, and others with $\Mass_\omega(u)<\int_X\omega^3$. Consequently $v_-(\omega)<\int_X\omega^3<v_+(\omega)$.
\end{enumerate}
\end{proposition}
\begin{proof}
(1) Since $\ddc|z|^2$ has constant coefficients and $y_j$ is pluriharmonic,
\[
 \ddc\beta=\ddc(y_1y_2)\wedge\ddc|z|^2,
\]
where
\[
 \ddc(y_1y_2)=\frac{\mathrm i}{2\pi}\bigl(\partial y_1\wedge\overline\partial y_2
       +\partial y_2\wedge\overline\partial y_1\bigr)
       =\frac{\mathrm i}{8\pi}\bigl(\mathrm dw_1\wedge\mathrm d\overline w_2
       +\mathrm dw_2\wedge\mathrm d\overline w_1\bigr).
\]
The Hermitian matrix of $\ddc(y_1y_2)$ in the basis $\mathrm dw_1,\mathrm dw_2$ is a multiple of $\left(\begin{smallmatrix}0&1\\1&0\end{smallmatrix}\right)$, which is indefinite; hence so is $\ddc\omega=\ddc\beta$. Since $\theta$ is closed and $\beta^2=0$, we have $\ddc(\omega^2)=2\theta\wedge\ddc\beta$, and this vanishes because every term of $\theta\wedge(\mathrm dw_1\wedge\mathrm d\overline w_2+\mathrm dw_2\wedge\mathrm d\overline w_1)$ repeats a differential.

(2) By \cref{thm:mass-expansion} with $n=3$ and $\ddc(\omega^2)=0$, $\Mass_\omega(u)-\int_X\omega^3=3\int_Xu\,\ddc u\wedge\ddc\omega$. The form $\ddc\omega$ is closed, so Stokes' theorem gives $\int_Xu\,\ddc u\wedge\ddc\omega=-\int_X\mathrm du\wedge\dc u\wedge\ddc\omega$. Now $\mathrm du\wedge\dc u=(\mathrm i/2\pi)\,\partial u\wedge\overline\partial u$, and a direct computation in the coordinates $w_1,w_2,z$ gives
\[
 \mathrm i\,\partial u\wedge\overline\partial u\wedge
 \mathrm i\bigl(\mathrm dw_1\wedge\mathrm d\overline w_2+\mathrm dw_2\wedge\mathrm d\overline w_1\bigr)
 \wedge\mathrm i\,\mathrm dz\wedge\mathrm d\overline z
 =-2\operatorname{Re}\Bigl(\frac{\partial u}{\partial w_1}\,
       \overline{\frac{\partial u}{\partial w_2}}\Bigr)\mathrm dV_0.
\]
Collecting the constants $(1/2\pi)(1/8\pi)(1/2\pi)$ gives the formula.

(3) Since $\ddc u$ has no components in the $z$-direction, $\omega+\ddc u=(\theta+\ddc u)+\beta$ splits into blocks, and $u\in\PSH(X,\omega)$ if and only if $u\in\PSH(X,\theta)$. By \cref{thm:ot-transverse-potentials}, $u=f(y)+\log(y_1y_2)$ with $f$ convex, so $g$ is locally Lipschitz, hence Lipschitz by periodicity. In the variables $\varsigma_j=\log y_j$, convexity of $f$ is equivalent to positivity of the symmetric matrix
\[
 H=I+D^2g-\operatorname{diag}(\partial_1g,\partial_2g),
\]
as one checks from $y_iy_j\,\partial^2f/\partial y_i\partial y_j=\partial^2\Phi/\partial\varsigma_i\partial\varsigma_j-\delta_{ij}\,\partial\Phi/\partial\varsigma_i$ for $\Phi(\varsigma)=f(\mathrm e^{\varsigma})=g(\varsigma)-\varsigma_1-\varsigma_2$. This condition is linear in $g$ with constant coefficients, so it is preserved by convolution with a smooth kernel on $\mathbb R^2$; the regularized functions are periodic, and they converge uniformly to $g$ with gradients converging almost everywhere. By the continuity of bounded Monge--Amp\`ere products under uniform convergence and dominated convergence, it suffices to prove all assertions for smooth $g$ with $H\geq0$.

For such $g$, the function $u$ depends only on $y$, so $\partial u/\partial w_j=-\frac{\mathrm i}2\partial u/\partial y_j$ and the integrand in (2) is $\frac14\partial_{y_1}u\,\partial_{y_2}u=\partial_1g\,\partial_2g/(4y_1y_2)$. A direct computation gives $\omega^3=3\theta^2\wedge\beta=\frac{3}{64\pi^3}\,\mathrm dV_0/(y_1y_2)$. By \eqref{eq:ot-fiber-integration} and (2),
\[
 \Mass_\omega(u)=\frac{\int_X\omega^3}{|T_U|}\int_{T_U}
       \bigl(1+\partial_1g\,\partial_2g\bigr)\,\mathrm d\varsigma
 =\frac{\int_X\omega^3}{|T_U|}\int_{T_U}PQ\,\mathrm d\varsigma,
\]
where the second equality uses $\int_{T_U}\partial_jg\,\mathrm d\varsigma=0$.

The diagonal entry $H_{11}=P-\partial_1P\geq0$ says that $\mathrm e^{-\varsigma_1}P$ is nonincreasing in $\varsigma_1$, that is,
\begin{equation}\label{eq:ot-p-monotone}
 P(\varsigma+\tau e_1)\leq\mathrm e^{\tau}P(\varsigma)\qquad(\tau\geq0).
\end{equation}
If $P(\varsigma)<0$, the right side tends to $-\infty$ as $\tau\to\infty$, which is impossible for the bounded function $P$; hence $P\geq0$, and likewise $Q\geq0$. Thus $\partial_jg\leq1$. The lattice $\Lambda_U$ contains two independent vectors $\lambda,\lambda'$ with positive coordinates, and every point of $T_U$ has the form $\varsigma+\tau\lambda+\tau'\lambda'$ with $\tau,\tau'\in[0,1)$. Along $\tau\mapsto\varsigma+\tau\lambda$, the derivative of $g$ is at most $|\lambda|_1=\lambda_1+\lambda_2$, and its integral over $[0,1]$ vanishes by periodicity, so $|g(\varsigma+\tau\lambda)-g(\varsigma)|\leq|\lambda|_1$. Hence $\operatorname{osc}g\leq L$, with $L=|\lambda|_1+|\lambda'|_1$ depending only on $\Lambda_U$.

Averaging \eqref{eq:ot-p-monotone} in the form $P(\varsigma-\tau e_1)\geq\mathrm e^{-\tau}P(\varsigma)$ over $\tau\in[0,1]$ gives
\[
 (1-\mathrm e^{-1})P(\varsigma)\leq\int_0^1P(\varsigma-\tau e_1)\,\mathrm d\tau
 =1-g(\varsigma)+g(\varsigma-e_1)\leq1+L.
\]
On the other hand, averaging $P(\varsigma)\geq\mathrm e^{-\tau}P(\varsigma+\tau e_1)$ over $\tau\in[0,R]$ with $R=\max(2L,1)$ gives
\[
 P(\varsigma)\geq\frac{\mathrm e^{-R}}R\int_0^RP(\varsigma+\tau e_1)\,\mathrm d\tau
 =\frac{\mathrm e^{-R}}R\bigl(R-g(\varsigma+Re_1)+g(\varsigma)\bigr)
 \geq\frac{\mathrm e^{-R}}2.
\]
The same bounds hold for $Q$. Therefore the average of $PQ$ over $T_U$ lies between $\mathrm e^{-2R}/4$ and $\bigl((1+L)/(1-\mathrm e^{-1})\bigr)^2$, which proves (3).

(4) Let $k$ belong to the dual lattice $\Lambda_U^*=\{k\in\mathbb R^2:\langle k,\lambda\rangle\in\mathbb Z\text{ for all }\lambda\in\Lambda_U\}$, and put $g_\epsilon(\varsigma)=\epsilon\cos(2\pi\langle k,\varsigma\rangle)$. For small $\epsilon>0$, the matrix $H$ of $g_\epsilon$ is positive definite, so $u_\epsilon=g_\epsilon(\varsigma)$ is a smooth function with $\omega+\ddc u_\epsilon>0$. By (3),
\[
 \frac{\Mass_\omega(u_\epsilon)}{\int_X\omega^3}
 =1+\frac1{|T_U|}\int_{T_U}\partial_1g_\epsilon\,\partial_2g_\epsilon\,\mathrm d\varsigma
 =1+2\pi^2\epsilon^2k_1k_2.
\]
Since the lattice $\Lambda_U^*$ has vectors in every open quadrant, we may choose $k$ with $k_1k_2>0$, respectively $k_1k_2<0$, which gives the two inequalities.
\end{proof}

For a general smooth potential, averaging over the fibers separates the invariant part from the fluctuations. For $u\in\PSH(X,\omega)$ let $Au$ be the average of $u$ over the fibers of $X\to T_U$, that is, over the translations $(w,z)\mapsto(w+r,z+\zeta)$ with $(r,\zeta)\in(\mathbb R^2\times\mathbb C)/(\sigma_{\mathbb R},\sigma_{\mathbb C})(\mathcal O_K)$. These translations preserve $\omega$ on the universal cover, are normalized by $\Gamma$, and act on the fiber torus by measure-preserving maps, because the units have norm one. Hence $Au$ is a well-defined $\omega$-plurisubharmonic function depending only on $y$, and $c_U^{-1}\int_X\omega^3\leq\Mass_\omega(Au)\leq c_U\int_X\omega^3$ by \cref{prop:ot-threefold-mass}. For smooth $u$, the integral of a $y$-derivative of $u-Au$ over a fiber vanishes, and so does the integral of an $x$-derivative. The cross terms between $Au$ and $u-Au$ in \cref{prop:ot-threefold-mass}(2) therefore integrate to zero, and
\begin{equation}\label{eq:ot-fluctuation}
 \Mass_\omega(u)=\Mass_\omega(Au)+\frac{3}{16\pi^3}\int_X
 \operatorname{Re}\Bigl(\frac{\partial v}{\partial w_1}\,
       \overline{\frac{\partial v}{\partial w_2}}\Bigr)\mathrm dV_0,
 \qquad v=u-Au.
\end{equation}
Since the extrema of the masses of bounded potentials are already attained by smooth ones (\cref{sec:gaud-volume-extrema}), $v_+(\omega)<+\infty$ if and only if the fluctuation term in \eqref{eq:ot-fluctuation} is bounded above on smooth $\omega$-plurisubharmonic functions; whether $v_-(\omega)>0$ likewise depends only on how negative this term can be relative to $\Mass_\omega(Au)$.
